
\chapter{Single Qubits}

\section{ The unit of information}

Almost all articles or news reports talking about quantum computing start off by explaining how qubits are different from classical bits. Vauge statements about superposition and entanglement try to impress the reader with a general idea of how quantum computers could be useful. This chapter applies the maths from previous chapters to properly describe superposition. Entanglement requires a later chapter to be introduced mathematically. Earlier, the coin example was introduced with the heads \& tails outcomes we can observe. A qubit has the corresponding measurement outcomes $\ket{0}$ \& $\ket{1}$, we can just replace $\ket{H}$ with $\ket{0}$ and $\ket{T}$ with $\ket{1}$. 

$$
\ket{0} = \begin{bmatrix} 1 \\ 0 \end{bmatrix} ; \ket{1} = \begin{bmatrix} 0 \\ 1 \end{bmatrix} 
$$

These two states are called the computational basis.Just as with the coins, our qubits can be in a superposition of $\ket{0}$ \& $\ket{1}$. We can have any probability of measuring $\ket{0}$ \& $\ket{1}$, so long as all the probabilities are real positive numbers which add to 1. 

\section{ The Bloch Sphere }

So far the linear algebra used makes for an efficient computational tool to describe what quantum states are and how they evolve. We can model how the qubits change using vectors and calculate probabilities with the inner product. One limitation of this Dirac notation is that it is difficult to visualise quantum states. A column vector is hard to think of intuitively.Thankfully, everything that has been used so far can easily be represented on a diagram. A sphere with radius of 1 can represent the quantum state of any qubit. This sphere is known as the Bloch sphere. Below, the quantum state $\ket{\psi}$ is represented on the Bloch Sphere. 

\begin{figure}[ht!]
    \centering
    \includegraphics[scale=0.5]{images/Bloch_Sphere.png}
    \caption{The Bloch Sphere: a geometric representation of the state space of a qubit.\protect\footnote{A useful tool for understanding the Bloch sphere can be found here: \href{https://javafxpert.github.io/grok-bloch/}{https://javafxpert.github.io/grok-bloch/}.}}
    \label{fig:bloch_sphere}
\end{figure}



At the top of the sphere  is the state $\ket{0}$, at the bottom is the state $\ket{1}$. Our state $\ket{\psi}$ is somewhere in between so it is a superposition of $\ket{0}$ \& $\ket{1}$. The closer $\ket{\psi}$  is to the top, the greater the probability of measuring $\ket{0}$, the same is true for the bottom with $\ket{1}$. The angle $\theta$ between the $\ket{0}$ \& $\ket{\psi}$ gives us these probabilities. 

Notice how $\ket{0}$ and $\ket{1}$ are separated by $180^{\circ} (\pi)$ on the sphere, but they are perpendicular, so the angle between them is actually $90 ^{\circ} (\pi/2)$. The angles on the Bloch sphere are twice as large as the angles for the qubits.

\textbf{The angle $\theta$ on the Bloch sphere gives the probabilities with angle $\theta/2$}

$$ P(\ket{0}) = |\braket{0|\psi}|^2 = (\cos{\theta/2})^2 $$
Similarly for the probability of $\ket{1}$

$$P(\ket{1}) = |\braket{0|\psi}|^2 = (\sin{\theta/2})^2$$

Up and down, the z-axis is where the $\ket{0}$ \& $\ket{1}$ are, for that reason the z-axis is known as the computational basis. In addition we have the x \& y axis which correspond to different states. 

The states located on the ends of the axes are known as eigenstates. For $Z$ these are $\ket{0}$ \& $\ket{1}$. The $X$ axis has eigenstates of $\ket{+}$ and $\ket{-}$, and the $Y$ axis has eigenstates $\ket{+i}$ and $\ket{-i}$. 

In quantum mechanics, eigenstates are defined by a particular operator. When the operator acts on its eigenstate the same state is returned, multiplied by an eigenvalue. The most famous example is the Schrödinger equation,

$$ \hat{H}\ket{\psi} = E\psi $$

Here the operator $\hat{H}$ is known as the Hamiltonian, which acts on the eigenstate $\ket{\psi}$ to return the energy $E$ associated with being in that state. Coming back to our qubits, we can use $Z$ as an example operator. If the qubit is prepared in the initial starting state $\ket{0}$, we can apply the gate $Z$ to $\ket{0}$. Since $\ket{0}$ is an eigenstate of $Z$, the resulting quantum state is simply $\ket{0}$ multiplied by 1 (the corresponding eigenvalue). Instead, if we started with $\ket{1}$, we would return $-1\ket{1}$. The eigenvalue of $-1$ cannot be measured, any more than the eigenvalue of $+1$ can be. 

\section{It's not just a phase!} \label{It's not just a phase!}

To show probabilities, a single angle is enough. So why do we need a sphere? There is another important parameter that describes a qubit: the phase. At last, our coins will not work to describe phase, there is no analogue in classical computing, nor in simple probability. For example, earlier the flip coin operator $F$ had this form:

$$
F =  \frac{1}{\sqrt{2}} \begin{bmatrix} 1 & 1 \\ 1 & -1 \end{bmatrix}
$$

In quantum computing, this operator is known as a Hadamard gate with symbol $H$. 
If we apply $H$ to $\ket{0}$ we get the superposition state 

$$
H\ket{0} = \frac{1}{\sqrt{2}} ( \ket{0} + \ket{1}) = \ket{+}
$$

Applying it to $\ket{1}$ gives us a different superposition state: 

$$
H\ket{1} = \frac{1}{\sqrt{2}} ( \ket{0} - \ket{1}) = \ket{-}
$$

We call this state $\ket{-}$ because of the minus sign between the $\ket{0}$ \& $\ket{1}$

The probability of getting $\ket{0}$ or $\ket{1}$ is the same for both. If we were to measure either state, we would get 0 half the time and 1 the other half the time. Since the measurement gives the same outcome you might think they are the same state. They are different states, and surprisingly there is an easy way for us to tell which one we have. 

What happens if we apply $H$ on these states? 

$$ H\ket{+} = H \frac{1}{\sqrt{2}} ( \ket{0} + \ket{1})$$

$$= \frac{1}{\sqrt{2}} (H\ket{0} + H\ket{1})$$

$$ = \frac{1}{\sqrt{2}}\left( \frac{1}{\sqrt{2}}(\ket{0} + \ket{1})+\frac{1}{\sqrt{2}}(\ket{0} - \ket{1}))  \right)$$

The $\ket{1}$'s cancel each other out and the $\ket{0}$ add together to give us

$$ = \ket{0} $$

So if we apply $H$ once to $\ket{0}$ we get $\ket{+}$ and then applying it again goes back to $\ket{0}$. This implies applying $H$ twice does nothing. $H^2 = I$.

So if we ever have $\ket{+}$ we apply $H$ and get $\ket{0}$ 100\% of the time. If we get $\ket{1}$ we must have had another state to begin with. If we had $\ket{-}$ would measure 1 100\% of the time. In this example, where is the phase? 

The phase is contained in the + or - sign between the $\ket{0}$ \& $\ket{1}$. 

The phase can be more than just $\pm 1$, it can be any complex number with magnitude 1. In general, any quantum state can have a phase $\phi$ 

$$ \alpha\ket{0} + \beta e^{i\phi}\ket{1} $$

With $ 0 \leq \phi < 2\pi $

We can use what we know from the Bloch Sphere to write $\alpha$ \& $\beta$ in terms of the angle $\theta$

$$ \cos{(\theta/2)}\ket{0} + \sin{(\theta/2)} e^{i\phi}\ket{1} $$

Remember Euler's formula (section 3.1.2) which allows us to convert from the exponential to the complex number

$$ e^{i \phi} = cos(\phi) + i sin(\phi) $$

If we choose $\phi = 0$ we get $e^0 = 1$ which gives us 

The phase is also shown on the Bloch sphere. In the above diagram the state $\ket{\psi}$ is described by two angles $\theta$ describing the probabilities of measuring 0 or 1 and $\phi$ describing the phase. The angle $\phi$ corresponds to a rotation by $\phi$ around the z-axis in the Bloch sphere.  

Phases can be used to encode information, such as with the quantum Fourier transform, or to mark out quantum states (as with Grover's algorithm). 

\textbf{Phases only matter when they are relative.}

 \section{Qubit gates }

As mentioned in Chapter 4, operators are what allow us to do things to quantum states. Below are a collection of the most common single qubit operators. 
\begin{table}[ht!]
\centering

\label{tab:single_qubit_gates}
\begin{tabular}{|m{1cm}| m{5cm} | m{3cm} | m{3cm} | m{2cm}|} 
\hline
 Gate & Description & Effect on $\ket{0}$ & Effect on $\ket{1}$ & Matrix \\ [0.5ex] 
\hline
 $I$ & Identity: do nothing & $\ket{0} \rightarrow  \ket{0} $ & $\ket{1} \rightarrow \ket{1}$ & $\begin{bmatrix} 1 & 0 \\ 0 & 1 \end{bmatrix}$ \\ 
\hline
 $X$ & Not: Swaps $\ket{0}$ with $\ket{1}$   &$\ket{0} \rightarrow \ket{1}$ &$\ket{1} \rightarrow  \ket{0}$ & $\begin{bmatrix} 0 & 1 \\ 1 & 0 \end{bmatrix}$ \\
\hline
 $Y$ & Y: Rotation by angle $\pi$ around the y-axis of the Bloch sphere & $\ket{0} \rightarrow  i\ket{1} $ & $\ket{1} \rightarrow -i\ket{0} $ & $\begin{bmatrix} 0 & -i \\ i & 0 \end{bmatrix}$ \\
\hline
 $Z$ & Z: Adds a phase of $\pi$ to $\ket{1}$&$\ket{0} \rightarrow \ket{0}$&$\ket{1} \rightarrow -\ket{1}$ &$\begin{bmatrix} 1 & 0 \\ 0 & -1 \end{bmatrix}$  \\
\hline
 $H$ & Hadamard: Creates the superposition state & $\ket{0} \rightarrow \frac{1}{\sqrt{2}} ( \ket{0} + \ket{1})$ & $\ket{1} \rightarrow \frac{1}{\sqrt{2}} ( \ket{0} - \ket{1})$ & $\frac{1}{\sqrt{2}}\begin{bmatrix} 1 & 1 \\ 1 & -1 \end{bmatrix} $\\
\hline
 $S$ & Adds a phase of $\pi/2$ to $\ket{1}$ &$\ket{0} \rightarrow  \ket{0} $ &$\ket{1} \rightarrow e^{i\pi/2}\ket{1}$  & $\begin{bmatrix} 1 & 0 \\ 0 & e^{i\pi/2} \end{bmatrix}$\\
\hline
 $T$ & Adds a phase of $\pi/4$ to $\ket{1}$ &  $\ket{0} \rightarrow  \ket{0} $& $\ket{1} \rightarrow e^{i\pi/4}\ket{1}$ & $\begin{bmatrix} 1 & 0 \\ 0 & e^{i\pi/4} \end{bmatrix}$\\ 
\hline
\end{tabular}
\caption{Common single-qubit gates and their action on the computational basis}
\end{table}

Sometimes it is easier to perform gates by knowing what they do to the $\ket{0}$, $\ket{1}$ states, sometimes it is easier to do the matrix multiplication. 

A more comprehensive list of gates \href{https://raw.githubusercontent.com/qiskit-community/qiskit-textbook/main/content/ch-states/supplements/single-gates-cheatsheet.pdf}{can be found here}

\section{Parametrised quantum gates }

In addition to these fixed quantum gates, there are quantum gates we can specify ourselves. These gates are described by a parameter we can control to change what the gate does. These gates are very similar to nodes in a neural network where we can change the weights (we will see a lot more about these in a chapter on quantum machine learning).For instance, we might want to get more than a 50\% chance of measuring an output. We can see on the Bloch sphere that if we rotate our quantum state by some angle towards $\ket{0}$ or $\ket{1}$ we can boost the probability of measuring either state.


The matrix that does this rotation is known as $R_y$.

$$
R_y(\theta) = \begin{bmatrix} \cos(\theta /2) & -\sin(\theta /2) \\ \sin(\theta /2) & \cos(\theta /2) \end{bmatrix}
$$

\textbf{The angle we give as input (on the Bloch Sphere) is $\theta$ the angle we perform the rotation is $\theta/2$.}

We also have rotations for the x \& y axis too. In fact, the X,Y gates are special cases of $R_x(\theta)$ \& $R_y(\theta)$ with $\theta = \pi$. 

$R_z(\phi)$ (also known as P) is also useful because it allows us to give any phase to a qubit. 

$$
R_z(\phi) = P(\phi) = \begin{bmatrix} 1 & 0 \\ 0 & e^{i\phi} \end{bmatrix} = \begin{bmatrix} e^{-i\phi/2} & 0 \\ 0 & e^{i\phi/2} \end{bmatrix}
$$

This gives a phase shift of $\phi$ between $\ket{0}$ \& $\ket{1}$. 


\section{ Compiling any single qubit gate }

Just as we can represent any single qubit state with two parameters, we can also represent any single qubit gate using parameters. 

$$
U(\theta, \phi, \lambda) = \begin{bmatrix} \cos(\theta/2) & -e^{i\lambda}\sin(\theta/2) \\ e^{i\phi}\sin(\theta/2) & e^{i(\phi + \lambda)}\cos(\theta/2) \end{bmatrix} 
$$

Just as we can reach any point on the Bloch Sphere with $\theta$ \& $\phi$ we can get any quantum gate using $U$. 

The set of all gates that can be made from $U$ are known as the special unitary group $SU(2)$.

Since quantum computers are physical objects, not mathematical abstractions, $U$ requires some mechanism by which it can be programmed into our qubit. In practice our $SU(2)$ matrices need to be mapped into a quantum control sequence defined by physical parameters. For qubits controlled by microwaves and/or lasers these physical parameters define the pulse that implements the operator on the qubit. To do this, an intermediate representation is defined allowing the user to input any quantum operator they want and then passing on to the control electronics a set of instructions.

Intermediate representations are effective because the linear algebra by which operators are defined allows $U$ to be decomposed into a product of matrices. The matrices are chosen such that each matrix can be implemented effectively by the control electronic. We call these matrices native gates. Depending on how a qubit is built their native gate set will be different. This is especially true for entangling operations. 

% \section{ Operators in Ket-Bra Form}

% There is a very elegant way to express operators using Dirac notation. If one 

\section{Expectation values revisited}

In section \ref{}

\section{Gates vs operators, computation and physics}

In the last chapter, operators were introduced as matrices that act on a quantum state. In this chapter, it seems quantum gates do the same thing. This is true, all quantum gates are operators, but not all operators are quantum gates. This book is written in the language of quantum information procesing, where physical processes are abstracted away into idealised qubits. Physics is a theory of information, what it is that can be known about the universe, and so the language used, the mathematical framework, is generally equivalent to our idealised computations. As always, there are applications for which small details can become more convenient to adopt the physics perspective. 

For instance, the creation operator $\hat{a}^\dagger$ creates a photon from the vacuum but can't be represented by a $SU(2)$ matrix. The momentum operator $\hat{p} = -i\hbar\frac{\partial}{\partial x}$ is a partial derivative existing in a continuous momentum Hilbert space of momentum states. Being able to map our physical world onto qubits, or moving from the physics perspective to the computational perspective, is the key skill in making use of quantum computers. Doing this mapping enables quantum simulation. It generally considered that quantum computers are more effective than classical computers in this mapping.  

\section{ Bonus: Beyond Bits }

For this course, we work with qubits- 2 level quantum systems. There is considerable research into using more than 2 levels for quantum computation. Such systems are called qudits: d-level quantum systems that we can use to process quantum information. 

Qudits are much less commonly used than qubits for a few reasons
\begin{itemize}
\item Having multiple levels is much more complicated. Think about the X gate, for a qudit, this gate becomes much more complicated

\item Qudits present even more engineering challenges than qubits. For instance, a qubit in a trapped ion consists of the ground level and the first excited state. Higher energy levels have smaller gaps in energy, so they are more difficult to control precisely. This also introduces more complex physical error models. 

\item Quantum error correction for qudits is much more complicated. With qubits $\ket{0}$ can only go to $\ket{1}$ but for a 3-level system there are 2 possible states $\ket{0}$ could go to. And that's just one qutrit (a qudit with d = 3). Correlations between all 3 states of two or more qudits means the different combinations of errors scale exponentially. This is also true with qubits, but at least they have a base of 2 rather than 3 or more.
\end{itemize}

\section{Chapter Exercises}


\begin{enumerate}
    \item Where on the sphere? We know the superposition state $\ket{+}$ has equal probabilities of $\ket{0}$ \& $\ket{1}$. Where on the Bloch sphere would this state be represented? What is the angle $\theta$ that would describe this state?
    \item Back to the start. Apply $H$ twice to $\ket{1}$ and verify the state we end up with is $\ket{1}$. 
    \item  The coin doesn't work. When we apply the Hadamard twice to a state, we get the same state back. Without doing any maths, consider throwing a coin. Why does throwing the coin twice not guarantee you get back the side you started with? 
    \item By considering the effect on $\ket{0}$ \& $ \ket{1}$, or by doing matrix multiplication, calculate the following states:
    
    \begin{enumerate}[a.]
        \item \( XZ\ket{1} \)
        \item \( XH\ket{+} \)
        \item \( XYZ\ket{0} \)
    \end{enumerate}
    \item By using $\theta = \pi/2$ apply $R_y(\pi/2)$ to the state $\ket{1}$. What gate does $R_y(\pi/2)$ do? 
    \item The X gate turns $\ket{0}$ to $\ket{1}$ and $\ket{1}$ to $\ket{0}$. By considering angles on the Bloch sphere, what is the angle $\theta$ such that $R_y(\theta) = Y$
    \item Which gate in the list above is represented by  $U(\theta, \phi, \lambda) = U(0,0,0)$ 
\end{enumerate}



\section{Chapter 5 Summary }

\begin{itemize}
    \item The state of a qubit can be represented on the Bloch sphere
    \item The BS shows the probability and phase of a qubit 
    \item The phase of a qubit is the $e^{i\theta}$ difference between the $\ket{0}$ \& $\ket{1}$ states
    \item There are many qubit gates we can use to perform operations on qubits
    \item Some qubit gates are rotations specified by a parameter which describes what the gate does
    
\end{itemize}